\subsection{Introduction}
\label{IV.20.0}

La mayoría de las nociones y resultados de los \textsection\textsection20 y 21 se vinculan directamente con el cap.~I, y apenas dependen de los caps.~II a IV, salvo por el uso ocasional de las nociones de profundidad y de anillo local regular (en \sref{IV.20.6}, \sref{IV.21.11}, \sref{IV.21.13} y \sref{IV.21.15}), del «teorema principal» de Zariski en \sref{IV.20.4} y \sref{IV.21.12}, y de las propiedades de las inmersiones transversalmente regulares en \sref{IV.20.6} y \sref{IV.21.15}.

\oldpage[IV-4]{226}
En el \textsection20 introducimos varias variantes de la noción de aplicación racional, estudiada ya en \sref[I]{section:I.7} desde un punto de vista aún bastante próximo al clásico y, por ello, poco adaptado al caso de preesquemas no necesariamente reducidos.
Las nociones y resultados del \textsection20 se utilizan para desarrollar en el \textsection21 (n\textsuperscript{os} \sref{IV.21.1} a \sref{IV.21.7}) la noción general de divisor y sus propiedades más elementales.
Esta noción resulta especialmente cómoda cuando los anillos locales de los preesquemas considerados son noetherianos e íntegramente cerrados, y sobre todo cuando son además \emph{factoriales} (\sref{IV.21.6} y \sref{IV.21.7}), debido a que en este último caso se identifica con la noción de \emph{ciclo de codimensión~$1$} (combinación lineal de subpreesquemas irreducibles de codimensión~$1$).
En \sref{IV.21.9} se determinan los divisores sobre un preesquema noetheriano de dimensión~$1$, aunque no necesariamente normal, lo cual es útil para diversas aplicaciones.
Los n\textsuperscript{os} \sref{IV.21.11} y \sref{IV.21.12} proporcionan dos teoremas importantes, debidos respectivamente a Auslander--Buchsbaum y Van der Waerden, relacionados con la noción de anillo factorial (los n\textsuperscript{os} \sref{IV.21.9}, \sref{IV.21.11} y \sref{IV.21.12} son independientes entre sí).
En los n\textsuperscript{os} \sref{IV.21.13} y \sref{IV.21.14}, también independientes de los tres anteriores, estudiamos una variante útil de la noción de anillo local factorial: la de anillo local \emph{parafactorial}, introducida en particular en~\cite{IV-41} al desarrollar teoremas de comparación entre el grupo de Picard de un preesquema proyectivo $X$ sobre un cuerpo $k$ y el de una «sección hiperplana».
Veremos en \sref{IV.21.14.1} (teorema de Ramanujam--Samuel) que los anillos locales parafactoriales son mucho más numerosos de lo que cabría esperar \emph{a priori}.

En \sref{IV.20.5}, \sref{IV.20.6} y \sref{IV.21.15} retomamos las nociones anteriores, pero desde un punto de vista «relativo» a un preesquema base fijo.
Por ahora estas nociones se utilizan con bastante poca frecuencia; en particular, la noción de divisor relativo apenas se emplea salvo cuando se trata de divisores positivos, y en ese caso puede explicitarse ventajosamente sin recurrir a la noción de función meromorfa relativa, mediante la noción de inmersión transversalmente regular de codimensión~$1$.
Por tanto, conviene omitir estas secciones en una primera lectura.
